% from ega0-16-fr.tex
\oldpage[0\textsubscript{IV-1}]{25}
Soient $A$ un anneau, $M$ un $A$-module de type fini, $\mathfrak{p}$ un idéal
premier de $A$. On a alors
\[
\label{0.16.1.8.1-fr}
  \dim_{A_{\mathfrak{p}}}(M_{\mathfrak{p}})
  +\dim_A(M/\mathfrak{p}M)\leq\dim_A(M).
\tag{16.1.8.1}
\]
\end{proposition}

\begin{proof}
On a
$\operatorname{Ann}(M_{\mathfrak{p}})=(\operatorname{Ann}(M))_{\mathfrak{p}}$,
donc les idéaux premiers de $A_{\mathfrak{p}}$ contenant
$\operatorname{Ann}(M_{\mathfrak{p}})$ correspondent biunivoquement aux
idéaux premiers de $A$ contenant $\operatorname{Ann}(M)$ et contenus dans
$\mathfrak{p}$. D'autre part, si $\operatorname{Ann}(M)=\mathfrak{n}$,
$V(\operatorname{Ann}(M/\mathfrak{p}M))=V(\mathfrak{p}+\mathfrak{n})$
($0_{\mathrm I}$, 1.7.5), donc les idéaux premiers de $A$ contenant
$\operatorname{Ann}(M/\mathfrak{p}M)$ contiennent
$\mathfrak{p}+\mathfrak{n}$~; d'où la conclusion.

Cette dernière remarque montre en outre que l'on a
\[
\label{0.16.1.8.2-fr}
  \dim_{A/\mathfrak{p}}(M/\mathfrak{p}M)=\dim_A(M/\mathfrak{p}M).
\tag{16.1.8.2}
\]
\end{proof}

\begin{proposition}[16.1.9]
\label{0.16.1.9-fr}
Soient $A$ un anneau noethérien, $\rho:A\to B$ un homomorphisme d'anneaux,
$M$ un $B$-module tel que $M_{[\rho]}$ soit un $A$-module de type fini. Alors
on a
\[
\label{0.16.1.9.1-fr}
  \dim_A(M_{[\rho]})=\dim_B(M).
\tag{16.1.9.1}
\]
\end{proposition}

\begin{proof}
En effet, l'anneau $B_M$ des homothéties du $B$-module $M$ s'identifie à un
sous-anneau de $C=\operatorname{End}_A(M_{[\rho]})$, et $C$ est un $A$-module
de type fini (Bourbaki, \emph{Alg.~comm.}, chap.~III, \S~3, no~1, lemme~2),
donc il en est de même de $B_M$~; d'ailleurs l'anneau $A_M$ des homothéties du
$A$-module $M_{[\rho]}$ est un sous-anneau de $B_M$, image canonique de $A$
dans $C$, donc $B_M$ est fini sur $A_M$~; la conclusion résulte donc de
(16.1.5).
\end{proof}

\begin{env}[16.1.10]
\label{0.16.1.10-fr}
Soit $A$ un anneau noethérien~; si $M$ est un $A$-module de longueur finie, on
sait (Bourbaki, \emph{Alg.~comm.}, chap.~IV, \S~2, no~5, prop.~7 et cor.~1 de
la prop.~7) que $\operatorname{Supp}(M)$ est un espace fini discret, donc
$\dim(M)=0$ si $M\neq0$~; réciproquement, si $M$ est un $A$-module de type
fini tel que $\dim(M)=0$, tout point de $\operatorname{Supp}(M)$ est fermé,
autrement dit est un idéal maximal de $A$, donc (\emph{loc.~cit.}, prop.~7)
$M$ est de longueur finie.
\end{env}

\subsection[Dimension d'un anneau semi-local noethérien]{Dimension d'un
anneau semi-local noethérien\protect\footnote{Ce numéro reproduit l'exposé
donné par Serre dans un cours au Collège de France en 1958.}}
\label{subsection:0.16.2-fr}

\begin{env}[16.2.1]
\label{0.16.2.1-fr}
Soient $A$ un anneau semi-local noethérien, $\mathfrak{r}$ son radical~;
rappelons qu'un idéal de définition $\mathfrak{q}$ de $A$ est un idéal tel que
$\mathfrak{q}\subset\mathfrak{r}$ et que $\mathfrak{q}$ contienne une
puissance de $\mathfrak{r}$~; il revient au même de dire que les seuls idéaux
premiers contenant $\mathfrak{q}$ sont maximaux, ou que $A/\mathfrak{q}$ est
un anneau artinien. Soit $M$ un $A$-module de type fini~; pour tout entier
$n>0$, $M/\mathfrak{q}^nM$ est un $A$-module de longueur finie, égale à
$\sum_{i=0}^{n-1}\operatorname{long}(\gr_i(M))$, où $\gr_\bullet(M)$ est le
$(A/\mathfrak{q})$-module gradué associé à $M$ muni de la filtration
$\mathfrak{q}$-préadique. On sait qu'il existe un polynôme $Q(n)$ à
coefficients rationnels tel que $\operatorname{long}(\gr_n(M))=Q(n)$ pour $n$
assez grand (III, 2.5.4)~; on en déduit aussitôt qu'il existe un polynôme
$P_{\mathfrak{q}}(M,n)$ en $n$, à coefficients rationnels, tel que
\[
  \operatorname{long}(M/\mathfrak{q}^nM)=P_{\mathfrak{q}}(M,n)
\]
pour $n$ assez grand~; il est clair que ce polynôme est unique et que son
coefficient dominant est $>0$ si $M\neq0$~; on dit que c'est le polynôme de
Hilbert-Samuel de $M$ pour la filtration $\mathfrak{q}$-préadique.
\end{env}

\begin{lemma}[16.2.2.1]
\label{0.16.2.2.1-fr}
